For our observational study, we adopt a causal inference framework to quantify the association between medication non-adherence and adverse event time.
Under the potential outcomes framework, these estimates approximate causal effects if unmeasured confounding is limited and positivity holds.
To help satisfy some of these assumptions, we apply preprocessing techniques such as trimming, detailed in Appendix~\ref{apd:preprocessing}.
In Appendix~\ref{adp:from_assoc_to_causal}, we outline the assumptions required for a valid causal interpretation, including those specific to survival analysis.

\paragraph{Restricted Mean Event Time (RMET).}  
RMET measures the expected time until an adverse event occurs, restricted to a predefined follow-up period.
Intuitively, it represents the average time a patient remains event-free within the observation window.\footnote{In survival analysis literature, this quantity is commonly referred to as the Restricted Mean Survival Time (RMST), but we adapt the terminology to emphasize its focus on adverse events.}
For patient \(i\), the RMET is defined as:
\[
\bar{T_i} = \mathbb{E}[\widetilde{T_i} | A_{i\tau}, \mathbf{Z}_{i\tau}] = \int_{0}^{M} {S}(u | A_{i\tau}, \mathbf{Z}_{i\tau}) \, du,
\]
where \( {S}(u | A_{i\tau}, \mathbf{Z}_{i\tau}) \) is the true survival function and \( M = 96 \) months denotes the maximum follow-up duration. In practice, we replace the true survival function with an estimated version using a survival model from Section~\ref{sec:methods_survival}.

\paragraph{Potential outcomes.} 
Under the potential outcomes framework, we denote the potential adverse event time under treatment \(a \in \{0,1\}\) by \(\widetilde{T_i}(a)\).
For potential outcomes, we can then define:
\[
\bar{T}_i(a) = \mu_a(\mathbf{Z}_{i\tau}) =\int_{0}^{M} S(u | A_{i\tau} = a, \mathbf{Z}_{i\tau})\, du,
\]
and the individual treatment effect (ITE) as:
\[
\text{ITE}_i = \bar{T}_i(1) - \bar{T}_i(0).
\]
By a standard derivation (see Appendix~\ref{apd:ate-ite}), the average treatment effect (ATE) is given by:
\(
\psi = \mathbb{E}\big[\bar{T}_i(1) - \bar{T}_i(0)\big].
\)
Finally, we can estimate the ATE empirically as:
\[
\widehat{\psi} = \frac{1}{N} \sum_{i=1}^{N}\left(\widehat{\mu}_1(\mathbf{Z}_{i\tau}) - \widehat{\mu}_0(\mathbf{Z}_{i\tau})\right),
\]
with
\(
\widehat{\mu}_a(\mathbf{Z}_{i\tau}) = \int_{0}^{M} \widehat{S}(u | A_{i\tau} = a, \mathbf{Z}_{i\tau})\, du.
\)

\paragraph{Causal estimators.}  
We extend meta-learning approaches (T-learner and S-learner from \citet{kunzel2019metalearners}) and matching methods to survival analysis, with algorithms provided in Appendix~\ref{apd:causal-meta-learner-algorithms}, introducing novel extensions for estimating treatment effects under censoring:
% Additionally, we employ causal survival forests as a nonparametric baseline.
\begin{itemize}[itemsep=0pt]
    \item \textbf{T-Learner:}
    Fit separate survival models for the treated (\(a=1\)) and control (\(a=0\)) groups. 
    The survival function is estimated independently for each treatment condition:
    \(
    \widehat{S}_a (u | \mathbf{Z}_{i\tau}) = {\widehat{S}(u | A_{i\tau} = a, \mathbf{Z}_{i\tau})},
    \)
    yielding two RMET estimators:
    \(
    \widehat{\mu}_a (\mathbf{Z}_{i\tau}) = \int_0^{M} \widehat{S}_a (u | \mathbf{Z}_{i\tau}) \, du,
    \)
    with \(a\in\{0,1\}\). The ATE is then estimated as:
    \[
    \widehat{\psi}_{\text{T-learner}} = \frac{1}{N} \sum_{i=1}^{N} \left( \widehat{\mu}_1 (\mathbf{Z}_{i\tau}) - \widehat{\mu}_0 (\mathbf{Z}_{i\tau}) \right).
    \]

    \item \textbf{S-Learner:} Fit a single survival model incorporating treatment as an additional covariate:
    \(
    \widehat{S}(u | \mathbf{Z}_{i\tau}, A_{i\tau}) = \widehat{S}(u | A_{i\tau}, \mathbf{Z}_{i\tau}).
    \)
    The single RMET estimator is then:
    \(
    \widehat{\mu} (\mathbf{Z}_{i\tau}, a) = \int_0^{M} \widehat{S}(u | \mathbf{Z}_{i\tau}, a) \, du,
    \)
    and the ATE is estimated by plugging \(a=0\text{ and } 1\) in the single estimator:
    \[
    \widehat{\psi}_{\text{S-learner}} = \frac{1}{N} \sum_{i=1}^{N} \left( \widehat{\mu} (\mathbf{Z}_{i\tau},1) - \widehat{\mu} (\mathbf{Z}_{i\tau},0) \right).
    \]

    \item \textbf{Matching (k-Nearest Neighbors)} \citep{stuart2010matching}\textbf{:}  
    Match individuals based on baseline covariates. 
    First set: \(
    \widehat{\mu}_{A_{i\tau}} (\mathbf{Z}_{i\tau}) = \widehat{\mu} (\mathbf{Z}_{i\tau}, A_{i\tau})
    \).
    Then for patient \(i\), let \(J_K(i)\) denote the set of \(K\) nearest neighbors from the opposite treatment group. The estimated counterfactual RMET is
    \[
    \widehat{\mu}_{1-A_{i\tau}}(\mathbf{Z}_{i\tau}) = \frac{1}{K} \sum_{j \in J_K(i)} \widehat{\mu}(\mathbf{Z}_{j\tau}, A_{j\tau}),
    \]
    and the ATE is estimated as:
    \footnotesize
    \[
    \widehat{\psi}_{\text{match}} = \frac{1}{N} \sum_{i=1}^{N} \left(\widehat{\mu}_{A_{i\tau}}(\mathbf{Z}_{i\tau}) - \widehat{\mu}_{1-A_{i\tau}}(\mathbf{Z}_{i\tau})\right)(2A_{i\tau} - 1).
    \]
    \normalsize
    We consider \(K \in \{1,5,20\}\).
    See Appendix~\ref{apd:matching_derivation} for a more detailed derivation.
    
    \item \textbf{Causal Survival Forest (CSF)} \citep{cui2023estimating}\textbf{:} A nonparametric method that estimates heterogeneous treatment effects in censored survival data:
    % CSF partitions the feature space into homogeneous subgroups and estimates ITEs within each partition:
    \[
    \widehat{\theta}_{\text{CSF}}(\mathbf{Z}_{i\tau}) = \widehat{\mu}_1 (\mathbf{Z}_{i\tau}) - \widehat{\mu}_0 (\mathbf{Z}_{i\tau}),
    \]
    with the ATE computed as:
    \[
    \widehat{\psi}_{\text{CSF}} = \frac{1}{N} \sum_{i=1}^{N} \widehat{\theta}_{\text{CSF}}(\mathbf{Z}_{i\tau}).
    \]
    (Appendix~\ref{apd:csf_ate_derivation} provides a more detailed derivation of CSF).
\end{itemize}